%% ----------------------------------------------------------------------------
%%  §6.5  Il teorema di L'Hôpital
%% ----------------------------------------------------------------------------
\section{Il teorema di L'Hôpital}
\label{sec:hopital}

Siano $f$ e $g$ due funzioni che tendono a zero per $x\to x_{0}$. Abbiamo visto nel
\cref{cap:funzioni-continue} che il rapporto $\frac{f(x)}{g(x)}$ è allora una forma
indeterminata: in generale non è possibile dedurre il risultato del limite dai limiti
di $f$ e di $g$, ma occorre trasformare preliminarmente il rapporto in modo da togliere
l'indeterminazione. Il teorema di L'Hôpital serve a questo scopo, e sostituisce al
rapporto delle funzioni il rapporto delle loro derivate.

L'idea si vede bene nella \cref{fig:hopital-tangenti}. Se $f(x_{0})=g(x_{0})=0$, vicino a
$x_{0}$ le due funzioni sono approssimate dalle loro rette tangenti,
$f(x)\approx f'(x_{0})(x-x_{0})$ e $g(x)\approx g'(x_{0})(x-x_{0})$; il fattore comune
$x-x_{0}$, responsabile dell'indeterminazione, si semplifica, e il rapporto
$\frac{f(x)}{g(x)}$ tende al rapporto delle pendenze $\frac{f'(x_{0})}{g'(x_{0})}$.

\begin{figure}[htb]
\centering
\begin{tikzpicture}
\begin{axis}[width=0.75\textwidth, height=6cm, xmin=-0.75, xmax=0.75, ymin=-1.15,
  ymax=1.15, xtick={-0.5,0.5}, ytick={-1,1}, xlabel={$x$}, ylabel={$y$}, clip=false]
  \addplot[curva, domain=-0.6:0.6, samples=100] {exp(x)-1};
  \addplot[curva2, domain=-0.6:0.6, samples=100] {sin(deg(2*x))};
  \addplot[tratteggio, domain=-0.6:0.6, samples=2] {x};
  \addplot[tratteggio, domain=-0.52:0.52, samples=2] {2*x};
  \node[font=\scriptsize, rame, anchor=west] at (axis cs:0.6,0.96) {$g(x)=\sen2x$};
  \node[font=\scriptsize, blu, anchor=west] at (axis cs:0.6,0.79) {$f(x)=e^{x}-1$};
  \node[font=\scriptsize, grigio, anchor=west] at (axis cs:0.6,0.6) {$y=x$};
  \node[font=\scriptsize, grigio, anchor=south] at (axis cs:0.52,1.04) {$y=2x$};
\end{axis}
\end{tikzpicture}
\caption{Vicino all'origine $f(x)=e^{x}-1$ e $g(x)=\sen2x$ si confondono con le loro
tangenti $y=x$ e $y=2x$ (tratteggiate): il rapporto $\frac{f(x)}{g(x)}$ si comporta come
$\frac{x}{2x}=\frac12$, il rapporto delle pendenze.}
\label{fig:hopital-tangenti}
\end{figure}

\begin{teorema}[di L'Hôpital]\label{teo:hopital}\index{teorema!di L'Hôpital}
Sia $I$ un intervallo aperto contenente il punto $x_{0}$ e siano $f$, $g$ funzioni
derivabili in $I\setminus\{x_{0}\}$, tali che
\begin{equation}\label{eq:hopital-ipotesi}
  \lim_{x\to x_{0}}f(x)=0, \qquad \lim_{x\to x_{0}}g(x)=0 .
\end{equation}
Se $g'(x)\neq0$ per ogni $x\in I\setminus\{x_{0}\}$ e se esiste, finito o infinito, il
limite del rapporto delle derivate, allora esiste anche il limite del rapporto delle
funzioni, e risulta
\begin{equation}\label{eq:hopital}
  \lim_{x\to x_{0}}\frac{f(x)}{g(x)}=\lim_{x\to x_{0}}\frac{f'(x)}{g'(x)} .
\end{equation}
Il teorema vale anche:
\begin{enumerate}[label=\textup{(\roman*)}]
  \item per i limiti destri ($x\to x_{0}^{+}$) e sinistri ($x\to x_{0}^{-}$), con $f$, $g$
  derivabili soltanto da una parte di $x_{0}$;
  \item per le forme indeterminate $\frac\infty\infty$, cioè supponendo, al posto della
  \eqref{eq:hopital-ipotesi}, che $f$ e $g$ tendano all'infinito (anzi, basta che
  $\abs{g(x)}\to+\infty$, senza alcuna ipotesi sul limite di $f$);
  \item per $x\to+\infty$ e per $x\to-\infty$, con $f$, $g$ derivabili in una semiretta.
\end{enumerate}
\end{teorema}

\begin{schema}
\Ipotesi{\begin{itemize}[nosep, leftmargin=1.2em]
  \item $f,g$ derivabili in $I\setminus\{x_{0}\}$, con $g'\neq0$;
  \item $f(x)\to0$ e $g(x)\to0$ (oppure $\abs{g(x)}\to+\infty$) per $x\to x_{0}$;
  \item esiste $\displaystyle\lim_{x\to x_{0}}\frac{f'(x)}{g'(x)}=\ell\in\overline{\R}$.
\end{itemize}}
\Tesi{$\displaystyle\lim_{x\to x_{0}}\frac{f(x)}{g(x)}=\ell$.}
\end{schema}

Le ipotesi implicano che anche $g(x)\neq0$ per $x\neq x_{0}$ vicino a $x_{0}$, cosicché
il rapporto $\frac{f}{g}$ è definito. La dimostrazione nel caso generale richiede un
nuovo strumento, il teorema di Cauchy, ed è svolta nel \cref{sec:cauchy-hopital} in
appendice. Qui ci limitiamo a provare il teorema in un caso particolare, ma
significativo, che rende rigorosa l'idea della \cref{fig:hopital-tangenti}.

\begin{proposizione}[Il caso delle funzioni derivabili in $x_{0}$]\label{prop:hopital-semplice}
Siano $f$, $g$ definite in un intorno di $x_{0}$, derivabili in $x_{0}$, con
\[
  f(x_{0})=g(x_{0})=0, \qquad g'(x_{0})\neq0 .
\]
Allora $g(x)\neq0$ per $x\neq x_{0}$ in un intorno di $x_{0}$, e
\begin{equation}\label{eq:hopital-semplice}
  \lim_{x\to x_{0}}\frac{f(x)}{g(x)}=\frac{f'(x_{0})}{g'(x_{0})} .
\end{equation}
Se inoltre $f'$ e $g'$ esistono in un intorno di $x_{0}$ e sono continue in $x_{0}$, il
secondo membro è uguale a $\lim_{x\to x_{0}}\frac{f'(x)}{g'(x)}$.
\end{proposizione}

\begin{schema}
\Ipotesi{$f,g$ derivabili in $x_{0}$; $f(x_{0})=g(x_{0})=0$; $g'(x_{0})\neq0$.}
\Tesi{$\displaystyle\frac{f(x)}{g(x)}\to\frac{f'(x_{0})}{g'(x_{0})}$ per $x\to x_{0}$.}
\end{schema}

\begin{strategia}
Poiché $f(x_{0})=g(x_{0})=0$, possiamo scrivere $f(x)=f(x)-f(x_{0})$ e
$g(x)=g(x)-g(x_{0})$; dividendo numeratore e denominatore per $x-x_{0}$, il rapporto
$\frac{f}{g}$ diventa il quoziente di due rapporti incrementali. Ciascuno ha limite
finito, e quello del denominatore è diverso da zero: si applica il teorema sul limite
del quoziente.
\end{strategia}

\begin{dimostrazione}
Per $x\neq x_{0}$ poniamo
\[
  F(x)=\frac{f(x)-f(x_{0})}{x-x_{0}}=\frac{f(x)}{x-x_{0}}, \qquad
  G(x)=\frac{g(x)-g(x_{0})}{x-x_{0}}=\frac{g(x)}{x-x_{0}} .
\]
Per definizione di derivata, $F(x)\to f'(x_{0})$ e $G(x)\to g'(x_{0})\neq0$. Per la
permanenza del segno (\cref{teo:confronto-funzioni}\,(i), applicato a $G$ oppure a $-G$),
esiste $\delta>0$ tale che $G(x)\neq0$, e quindi $g(x)\neq0$, per
$0<\abs{x-x_{0}}<\delta$. In tali punti
\[
  \frac{f(x)}{g(x)}=\frac{f(x)/(x-x_{0})}{g(x)/(x-x_{0})}=\frac{F(x)}{G(x)}
  \ \longrightarrow\ \frac{f'(x_{0})}{g'(x_{0})}
\]
per il teorema sulle operazioni con i limiti (\cref{teo:operazioni-funzioni}). Se infine
$f'$, $g'$ sono continue in $x_{0}$, risulta $\frac{f'(x)}{g'(x)}\to\frac{f'(x_{0})}{g'(x_{0})}$,
di nuovo per il limite del quoziente.
\end{dimostrazione}

\subsection{Esempi: le forme \texorpdfstring{$\frac00$}{0/0} e \texorpdfstring{$\frac\infty\infty$}{inf/inf}}

Calcoliamo il limite
\begin{equation}\label{eq:hopital-es1}
  \lim_{x\to0}\frac{e^{x}-1}{\sen2x} .
\end{equation}
Si tratta di una forma indeterminata $\frac00$, che verifica le ipotesi della
\cref{prop:hopital-semplice}: numeratore e denominatore sono nulli e derivabili in $0$, e
$g'(0)=2\cos0=2\neq0$. Quindi
\[
  \lim_{x\to0}\frac{e^{x}-1}{\sen2x}=\frac{e^{0}}{2\cos0}=\frac12 ,
\]
come suggeriva la \cref{fig:hopital-tangenti}. Un'altra situazione dello stesso tipo è
\[
  \lim_{x\to1}\frac{x^{2}-1}{\log x}=\frac{2x}{1/x}\Big|_{x=1}=2 .
\]

Consideriamo ora due limiti notevoli, analoghi ai limiti di successioni della
gerarchia degli infiniti (\cref{teo:gerarchia}), che si calcolano facilmente con il
teorema di L'Hôpital nella forma $\frac\infty\infty$. Sia $b$ un numero reale
positivo. Si ha
\begin{equation}\label{eq:hopital-log-potenza}
  \lim_{x\to+\infty}\frac{\log x}{x^{b}}=\lim_{x\to+\infty}\frac{1/x}{b\,x^{b-1}}
  =\lim_{x\to+\infty}\frac1{b\,x^{b}}=0 :
\end{equation}
il logaritmo tende all'infinito più lentamente di qualunque potenza positiva. Per il
confronto tra potenza ed esponenziale, sia $n$ il più piccolo intero $\ge b$ (cioè
$n=b$ se $b$ è intero, $n=\lfloor b\rfloor+1$ altrimenti). Derivando $n$ volte il
numeratore si
ottiene
\begin{equation}\label{eq:hopital-potenza-exp}
  \lim_{x\to+\infty}\frac{x^{b}}{e^{x}}=\lim_{x\to+\infty}\frac{b\,x^{b-1}}{e^{x}}=\dots
  =\lim_{x\to+\infty}\frac{b(b-1)\cdots(b-n+1)\,x^{b-n}}{e^{x}}=0 ,
\end{equation}
perché $b-n\le0$ e quindi $x^{b-n}$ è limitata per $x\ge1$. Il ragionamento va letto
\emph{da destra verso sinistra}: l'ultimo limite esiste ed è nullo; le forme intermedie
$\frac{x^{b-k}}{e^{x}}$ con $k<n$ sono del tipo $\frac\infty\infty$ (perché $b-k>0$), e il
teorema di L'Hôpital, applicato $n$ volte risalendo la catena, assicura che anche i
limiti precedenti esistono e valgono $0$. Le \eqref{eq:hopital-log-potenza} e
\eqref{eq:hopital-potenza-exp} estendono alle funzioni la gerarchia
$\log n\ll n^{b}\ll a^{n}$ delle successioni.

\subsection{Le altre forme indeterminate}

Il teorema di L'Hôpital è utile anche per le forme indeterminate $0\cdot\infty$ e
$\infty-\infty$, che si riconducono a un quoziente con semplici manipolazioni. Per il
prodotto si scrive $fg=\frac{f}{1/g}$, oppure $fg=\frac{g}{1/f}$, ottenendo una forma
$\frac00$ oppure $\frac\infty\infty$. Ad esempio, se $b>0$,
\begin{equation}\label{eq:hopital-xb-log}
  \lim_{x\to0^{+}}x^{b}\log x=\lim_{x\to0^{+}}\frac{\log x}{x^{-b}}
  =\lim_{x\to0^{+}}\frac{1/x}{-b\,x^{-b-1}}=\lim_{x\to0^{+}}\frac{x^{b}}{-b}=0 .
\end{equation}
Per una differenza si cerca di scriverla come un unico quoziente, ad esempio
riducendo allo stesso denominatore:
\begin{align*}
  \lim_{x\to0}\Bigl(\frac1{\sen x}-\frac1x\Bigr)&=\lim_{x\to0}\frac{x-\sen x}{x\sen x}
  =\lim_{x\to0}\frac{1-\cos x}{\sen x+x\cos x}\\
  &=\lim_{x\to0}\frac{\sen x}{2\cos x-x\sen x}=\frac02=0,
\end{align*}
dove il teorema è stato applicato due volte (le prime due forme sono $\frac00$, l'ultimo
rapporto è continuo in $0$).

Le forme indeterminate esponenziali $0^{0}$, $1^{\infty}$, $\infty^{0}$ si trattano
scrivendo $f(x)^{g(x)}=e^{g(x)\log f(x)}$ (per $f>0$): l'esponente $g(x)\log f(x)$ è una
forma $0\cdot\infty$, e per la continuità della funzione esponenziale
(\cref{teo:continuita-esponenziale,cor:limite-composta-continua}) basta calcolarne il
limite. Ad esempio, usando la \eqref{eq:hopital-xb-log} con $b=1$,
\begin{equation}\label{eq:hopital-xx}
  \lim_{x\to0^{+}}x^{x}=\lim_{x\to0^{+}}e^{x\log x}=e^{\,\lim_{x\to0^{+}}x\log x}=e^{0}=1 ;
\end{equation}
per la forma $1^{\infty}$,
\begin{equation}\label{eq:hopital-1-infinito}
  \lim_{x\to0}(1+\sen x)^{1/x}=e^{\,\lim\limits_{x\to0}\frac{\log(1+\sen x)}{x}}
  =e^{\,\lim\limits_{x\to0}\frac{\cos x}{1+\sen x}}=e ;
\end{equation}
e per la forma $\infty^{0}$, per la \eqref{eq:hopital-log-potenza} con $b=1$,
\[
  \lim_{x\to+\infty}x^{1/x}=e^{\,\lim\limits_{x\to+\infty}\frac{\log x}{x}}=e^{0}=1 .
\]

Il programma seguente verifica tutti questi limiti con la funzione \texttt{limit} di
\texttt{sympy}. Si osservi che \texttt{sympy} non usa il teorema di L'Hôpital, ma
algoritmi basati sugli sviluppi asintotici, di cui la formula di Taylor del
\cref{sec:taylor} è il primo esempio.

\begin{python}[Verifica dei limiti con sympy]
import sympy as sp

x = sp.symbols("x", real=True)
b = sp.Rational(1, 2)                    # un esponente positivo
esempi = [((sp.exp(x) - 1) / sp.sin(2*x), 0, "+-"),
          ((x**2 - 1) / sp.log(x), 1, "+-"),
          (sp.log(x) / x**b, sp.oo, "-"),
          (x**5 / sp.exp(x), sp.oo, "-"),
          (x**b * sp.log(x), 0, "+"),
          (1 / sp.sin(x) - 1 / x, 0, "+-"),
          (x**x, 0, "+"),
          ((1 + sp.sin(x))**(1/x), 0, "+-"),
          (x**(1/x), sp.oo, "-")]
for f, x0, verso in esempi:
    print(f"lim {str(f):>24} per x -> {x0}:",
          sp.limit(f, x, x0, verso))
\end{python}
\begin{risultato}
lim    (exp(x) - 1)/sin(2*x) per x -> 0: 1/2
lim        (x**2 - 1)/log(x) per x -> 1: 2
lim           log(x)/sqrt(x) per x -> oo: 0
lim             x**5*exp(-x) per x -> oo: 0
lim           sqrt(x)*log(x) per x -> 0: 0
lim           1/sin(x) - 1/x per x -> 0: 0
lim                     x**x per x -> 0: 1
lim      (sin(x) + 1)**(1/x) per x -> 0: E
lim                 x**(1/x) per x -> oo: 1
\end{risultato}

\subsection{Tre avvertenze}

\begin{attenzione}[Prima di derivare, controllare che la forma sia indeterminata]
Il teorema si applica soltanto alle forme $\frac00$ e ai rapporti il cui denominatore
tende all'infinito (in particolare alle forme $\frac\infty\infty$). Applicato a
sproposito dà risultati sbagliati: $\lim_{x\to0}\frac{1+x}{2+x}=\frac12$ per
sostituzione, mentre il rapporto delle derivate vale costantemente $1$. Inoltre si
derivano numeratore e denominatore \emph{separatamente}: la derivata del quoziente
$\bigl(\frac fg\bigr)'$ non ha nulla a che fare con il teorema.
\end{attenzione}

\begin{attenzione}[Il teorema fornisce una condizione sufficiente]
Se il limite di $\frac{f'}{g'}$ non esiste, il teorema non dice nulla: il limite di
$\frac fg$ può esistere ugualmente. Per $f(x)=x+\sen x$ e $g(x)=x$, per $x\to+\infty$,
\[
  \frac{f(x)}{g(x)}=1+\frac{\sen x}{x}\ \to\ 1, \qquad\text{mentre}\qquad
  \frac{f'(x)}{g'(x)}=1+\cos x
\]
non ha limite, perché oscilla tra $0$ e $2$ (\cref{cor:non-esistenza}). Il programma
seguente lo mostra numericamente nei punti $x=\frac{k\pi}{2}$. Un esempio dello stesso
tipo, con $x\to0$, è $\frac{x^{2}\sen(1/x)}{\sen x}$, che tende a $0$, mentre il rapporto
delle derivate contiene il termine $\frac{\cos(1/x)}{\cos x}$, privo di limite (si
ricordino le funzioni $x^{n}\sen\frac1x$ del \cref{sec:retta-tangente}).
\end{attenzione}

\begin{python}[Quando il rapporto delle derivate non ha limite]
import math

# f(x) = x + sen x, g(x) = x per x -> +oo: f/g -> 1, f'/g' oscilla
print("     x = k pi/2        f/g     f'/g' = 1 + cos x")
for k in (2000, 2001, 2002, 2003, 200000, 200001, 200002):
    x = k * math.pi / 2
    print(f"{x:15.4f}   {(x + math.sin(x)) / x:.6f}   "
          f"{1 + math.cos(x):.6f}")
\end{python}
\begin{risultato}
     x = k pi/2        f/g     f'/g' = 1 + cos x
      3141.5927   1.000000   2.000000
      3143.1634   1.000318   1.000000
      3144.7342   1.000000   0.000000
      3146.3050   0.999682   1.000000
    314159.2654   1.000000   2.000000
    314160.8362   1.000003   1.000000
    314162.4070   1.000000   0.000000
\end{risultato}

\begin{attenzione}[Il teorema di L'Hôpital non dimostra i limiti notevoli]
Si è tentati di «dimostrare» il limite notevole $\frac{\sen x}{x}\to1$ con il teorema di
L'Hôpital: $\lim\frac{\sen x}{x}=\lim\frac{\cos x}{1}=1$. Il calcolo è corretto, ma non è
una dimostrazione: la formula $D\sen x=\cos x$ è stata ricavata nel
\cref{sec:derivate-elementari} proprio dal limite $\frac{\sen x}x\to1$, e usarla per
dimostrare quel limite sarebbe un circolo vizioso. Lo stesso vale per
$\frac{e^{x}-1}{x}\to1$ e $\frac{\log(1+x)}{x}\to1$, da cui sono state ricavate le
derivate dell'esponenziale e del logaritmo. Il teorema di L'Hôpital è uno strumento per
calcolare limiti \emph{nuovi} a partire dalle derivate \emph{già note}.
\end{attenzione}

\begin{notastorica}[Il marchese e il suo maestro]
Guillaume François Antoine de L'Hôpital (1661--1704), marchese di Sainte-Mesme,
pubblicò nel 1696 l'\emph{Analyse des infiniment petits pour l'intelligence des lignes
courbes}, il primo manuale di calcolo differenziale della storia, che ebbe un enorme
successo e diffuse in Europa i metodi di Leibniz. La regola che porta il suo nome vi
compare nella sezione IX. Già allora si sapeva che il marchese aveva preso lezioni dal
giovane Johann Bernoulli, ma solo nel Novecento, con la pubblicazione della loro
corrispondenza, si scoprì che nel 1694 L'Hôpital aveva offerto a Bernoulli un vitalizio
in cambio della comunicazione esclusiva delle sue scoperte matematiche, con il diritto di
usarle a suo piacimento. Il teorema «di L'Hôpital» è dunque, con ogni probabilità, di
Bernoulli. Va detto, a onore del marchese, che nella prefazione dell'\emph{Analyse} egli
riconosceva il suo debito verso Leibniz e verso i fratelli Bernoulli. La grafia
originale del nome era «L'Hospital»; nella grafia francese moderna la \emph{s} muta è
sostituita dall'accento circonflesso.
\end{notastorica}
